\documentclass{report}
\usepackage{amsmath,amssymb}
\setlength{\parindent}{0mm}
\setlength{\parskip}{1em}
\begin{document}
\begin{center}
\rule{6in}{1pt} \
{\large Yeonjong Shin \\
{\bf Optimal least squares method for polynomial regression}}

Department of Mathematics \\ University of Utah \\ 155S 1400E Rm 233 \\ Salt Lake City \\ Utah 84112 USA
\\
{\tt shin@math.utah.edu}\\
Dongbin Xiu\end{center}

We present a quasi-optimal sample set for least squares regression. The
quasi-optimal set is designed in such a way that, for a given number of
samples, it can deliver the regression result as close as possible to the
result obtained by a (much) larger set of candidate samples. The
quasi-optimal set is determined by maximizing a quantity measuring the
mutual column orthogonality and the determinant of the model matrix. This
procedure is non-adaptive, in the sense that it does not depend on the
sample data. This is useful in practice, as it allows one to determine,
prior to the potentially expensive data collection
procedure, where to sample the underlying system. In addition to
presenting the theoretical motivation of the quasi-optimal set, we also
present its efficient implementation via a greedy algorithm, along with
several numerical examples to demonstrate its efficacy. Since the
quasi-optimal set allows one to obtain a near optimal regression result
under any affordable number of samples, it notably outperforms other
standard choices of sample sets.


\end{document}
